\documentclass{ohm_project_description}
\usepackage[english]{babel}

\projecttitle{Case Study on the Use of New Technologies in Companies}
\projectauthor{TTZ Nürnberger Land / TH Nürnberg}
\projectdate{2026}

\begin{document}

\maketitle
\thispagestyle{fancy}

The aim of this case study is to survey primarily small and medium-sized enterprises (SMEs) in the Nürnberger Land region regarding their use of new technologies such as artificial intelligence, mobile and stationary robotic systems, and the potential for automating tasks and processes.

This includes identifying technologies already in use, assessing willingness to adopt new technologies in production, quality management and technical communication, and understanding the challenges and barriers involved. Surveys are to be methodically prepared, suitable companies identified, and decision-makers interviewed and evaluated. The result is a summarising presentation of the findings and a brief assessment of potentially relevant technologies for this sector.

The \textbf{Technology Transfer Centre Nürnberger Land} aims to advance knowledge transfer from research into primarily small and medium-sized enterprises in the Nürnberger Land region, maintaining a broad network of companies in this region.

\section*{Work Packages}
\begin{itemize}[leftmargin=0.5cm]
    \setlength\itemsep{.1em}
    \item Methodical preparation and execution of company surveys
    \item Identification of suitable companies in the Nürnberger Land region
    \item Conducting interviews with decision-makers
    \item Evaluation and summarising presentation of the results
    \item Classification of relevant technologies (AI, robotics, automation)
\end{itemize}

\section*{Requirements}
\begin{itemize}[leftmargin=0.5cm]
    \setlength\itemsep{.1em}
    \item Interest in case studies and empirical research
    \item Broad technical background
    \item Communication skills and organisational ability
    \item Independent and goal-oriented working style
\end{itemize}

\vspace{0.5cm}
This topic is designed as a \textbf{project thesis}.

\vfill
\textcolor{ohm_red}{\rule{\linewidth}{0.4mm}}
\textbf{\textcolor{ohm_red}{TTZ Nürnberger Land / Mobile Robotics Lab}} \\
\begin{tabular}{@{}ll}
\textbf{Supervisor:} & Prof. Dr. Christian Pfitzner \\
\textbf{E-Mail:}     & \href{mailto:christian.pfitzner@th-nuernberg.de}{christian.pfitzner@th-nuernberg.de} \\
\end{tabular}

\end{document}
